\section{From depth two to scaling laws}
\label{sec:scaling}

The preceding section proves the semantic depth theorem: in the fixed cubical sealed-extension
calculus, binary public trace is necessary, and higher sealing-generated structural trace is
computed by cubical horn filling.  This section asks what that theorem implies for the size of
public signatures as sealed layers accumulate.

The answer is intentionally conditional.  The recurrence below is not a law of ordinary
library growth.  It is the arithmetic shadow of three prior choices:
\begin{enumerate}
\item a genuine chronological window has already been identified;
\item a new sealed layer is advertised as globally acting on the active interface in that
window; and
\item the bundled per-site convention is used, so one active primitive basis site contributes
one primitive public interaction schema for the whole payload package.
\end{enumerate}
Without these assumptions, one gets only a sparse footprint formula.  A local datatype, a
transparent helper definition, or an orthogonal library construction may touch only a few
active sites, or none at all.  Such declarations are not expected to satisfy a Fibonacci law.

A useful way to read the section is therefore:
\[
  \text{depth theorem}
  \quad+\quad
  \text{chronological factorization}
  \quad+\quad
  \text{full coupling}
  \quad\Longrightarrow\quad
  \text{affine recurrence}.
\]
The Fibonacci statement is only the constant-payload endpoint of this recurrence.  It is not
the conceptual center of the paper.  The conceptual center remains the depth-two theorem: the
reason the recurrence has two previous terms is that the fixed cubical calculus turns higher
structural trace into derived horn computation and then factors surviving primitive binary
support through the two most recent normalized signatures.

\subsection{Layerwise quantities and the active window}
\label{subsec:layerwise-quantities}

We begin by naming the quantities that enter the count.  All of them are computed after
normalization.  In particular, a different record packaging, a different bracketing of a
Sigma-telescope, or an explicit surface field that is later replaced by a cubical horn term
does not change the numbers below.

\begin{definition}[Layerwise payload, trace, and basis size]
\label{def:layerwise-payload-trace}
Let
\[
  e_n
\]
be the admissible surface declaration whose sealing exports the stage-
\(n\) layer \(L_n\), and let \(X_n := X(e_n)\) be its elaborated candidate in the raw calculus.
The \emph{payload contribution} of layer \(n\) is
\[
  \kappa_n := \kappa(X_n) := |S_{\mathrm{core}}(L_n)|.
\]
The \emph{primitive trace contribution} of layer \(n\) is
\[
  \mu_n := \mu(e_n) := |S_{\mathrm{str}}(L_n)|.
\]
The normalized active basis exported by the layer is the disjoint union of its payload and
primitive trace basis:
\[
  S_{\mathrm{bas}}(L_n)
  := S_{\mathrm{core}}(L_n) \sqcup S_{\mathrm{str}}(L_n),
\]
and hence
\[
  |S_{\mathrm{bas}}(L_n)| = \kappa_n + \mu_n.
\]
\end{definition}

The notation \(S_{\mathrm{bas}}(L_n)\) is deliberately basis-level notation.  It does not count
all textual declarations exported by a layer.  It counts irreducible active primitive sites
after opacity, normalization, and derived-field elimination.  If a large bundle is treated as
one active basis site by the chosen convention, then it contributes one site.  If the bundle is
split into several independent active sites in the counting normal form, then it contributes
several.  This is why the recurrence is a theorem about normalized public signatures rather
than source files.

\begin{definition}[Depth-\(d\) historical basis window]
\label{def:historical-basis-window}
Assume that a foundation or extension discipline has a chronological window of size \(d\).
At stage \(n+1\), the depth-\(d\) active historical basis is the tagged coproduct
\[
  I^{(d)}_{n,\mathrm{bas}}
  :=
  \coprod_{j=0}^{d-1} S_{\mathrm{bas}}(L_{n-j}).
\]
For bootstrap stages, where some layer \(L_{n-j}\) does not yet exist, the corresponding
summand is defined to be the empty basis.
\end{definition}

The tags matter.  The same primitive shape appearing in two different sealed layers is counted
as two historical basis sites, because a new globally acting layer must specify how it
interacts with each opaque export as an export.  Conversely, two syntactic names in one layer
that normalize to the same active basis site are counted once.

\begin{remark}[Indexing convention]
\label{rem:scaling-indexing}
The layer \(L_n\) is exported by the declaration \(e_n\).  Its payload contribution is
\(\kappa_n\), its primitive public trace contribution is \(\mu_n\), and its normalized active
basis size is \(\kappa_n+\mu_n\).  The declaration exporting \(L_{n+1}\) is counted against
the active window ending at \(L_n\).  Thus, in the depth-two cubical case, the recurrence for
\(\mu_{n+1}\) uses \(L_n\) and \(L_{n-1}\):
\[
  \mu_{n+1}
  = \mu_n + \mu_{n-1} + \kappa_n + \kappa_{n-1}.
\]
This equality is asserted after the two-layer window has been populated.  Earlier bootstrap
steps use the empty-interface convention from Definition~\ref{def:historical-basis-window}.
\end{remark}

The next distinction is the one that prevents the recurrence theorem from being read too
broadly.

\begin{definition}[Sparse footprint and full coupling]
\label{def:sparse-footprint-full-coupling}
A stage-\(n+1\) declaration has a \emph{sparse active footprint} of depth \(d\) if it selects a
finite sub-basis
\[
  \mathrm{Foot}^{(d)}_n \subseteq I^{(d)}_{n,\mathrm{bas}}
\]
of the historical basis window.  Its sparse trace cost is the size of that selected footprint.

The declaration is \emph{fully coupled} against the depth-\(d\) window when
\[
  \mathrm{Foot}^{(d)}_{n,\mathrm{full}} = I^{(d)}_{n,\mathrm{bas}}.
\]
In words: the new sealed layer is advertised as globally acting on every active primitive basis
site in the relevant public window.
\end{definition}

Sparse footprints are the default for ordinary extension.  A definition may use only the
list type, only one universe level, or only a small set of operations.  Full coupling is a
special foundational-core regime.  A modality, universe-level operation, or promoted structural
package may advertise an action on the entire active interface.  The recurrence theorem is
about this special endpoint.

\subsection{Sparse footprints before the full envelope}
\label{subsec:sparse-footprints}

The arithmetic begins with a finite-set decomposition.  Since the historical window is a
tagged coproduct of layer bases, any selected footprint decomposes uniquely into the part it
selects from each recent layer.

\begin{lemma}[Window cardinality]
\label{lem:window-cardinality}
Let a stage-\(n+1\) declaration use an active footprint
\[
  \mathrm{Foot}^{(d)}_n \subseteq I^{(d)}_{n,\mathrm{bas}}.
\]
Then
\[
  |\mathrm{Foot}^{(d)}_n|
  =
  \sum_{j=0}^{d-1}
  \bigl|\mathrm{Foot}^{(d)}_n \cap S_{\mathrm{bas}}(L_{n-j})\bigr|.
\]
If the declaration is fully coupled, then
\[
  |\mathrm{Foot}^{(d)}_{n,\mathrm{full}}|
  =
  \sum_{j=0}^{d-1}
  (\kappa_{n-j}+\mu_{n-j}).
\]
\end{lemma}

\begin{proof}
The historical basis is the tagged coproduct
\[
  I^{(d)}_{n,\mathrm{bas}}
  =
  \coprod_{j=0}^{d-1} S_{\mathrm{bas}}(L_{n-j}).
\]
A finite subset of a tagged coproduct is the disjoint union of its intersections with the
tagged summands.  Taking cardinalities gives the sparse formula.  In the fully coupled case,
each intersection is the entire summand.  Definition~\ref{def:layerwise-payload-trace} then
identifies the summand size with \(\kappa_{n-j}+\mu_{n-j}\).
\end{proof}

For a concrete depth-two picture, suppose the recent active interface has five primitive basis
sites:
\[
  A, B, C, D, E.
\]
A sparse declaration might touch only \(A\) and \(C\), while a fully coupled foundational-core
extension touches all five.

\begin{center}
\begin{tabular}{lccccc l}
\hline
Regime & \(A\) & \(B\) & \(C\) & \(D\) & \(E\) & selected sites \\
\hline
sparse depth 2 & yes & no & yes & no & no & \(\{A,C\}\) \\
fully coupled depth 2 & yes & yes & yes & yes & yes & \(\{A,B,C,D,E\}\) \\
\hline
\end{tabular}
\end{center}

The recurrence is obtained by reading the second row through the identity
\(|S_{\mathrm{bas}}(L_i)|=\kappa_i+\mu_i\).  The first row gives a sparse cost, but not the
full-coupling recurrence.

\subsection{The full-coupling affine recurrence}
\label{subsec:full-coupling-recurrence}

We can now state the general recurrence.  It is parameterized by a chronological window size
\(d\).  In the fixed cubical calculus, Section~5 supplies \(d=2\).  In a UIP-like collapse,
one may have \(d=1\).  Other window sizes are included only to make clear which part of the
proof is arithmetic and which part is semantic.

\begin{theorem}[Full-coupling affine recurrence]
\label{thm:full-coupling-affine-recurrence}
Let a sealed extension discipline have obligation depth realized by a chronological window of
size \(d\).  Assume that the stage-\(n+1\) declaration is fully coupled against the whole
active window:
\[
  \mathrm{Foot}^{(d)}_{n,\mathrm{full}} = I^{(d)}_{n,\mathrm{bas}}.
\]
Then the primitive trace cost of the next sealed layer satisfies
\[
  \mu_{n+1}
  =
  \sum_{j=0}^{d-1} (\mu_{n-j}+\kappa_{n-j})
  \qquad (n\geq d),
\]
with bootstrap stages interpreted using empty earlier layers.
\end{theorem}

\begin{proof}
By the trace-cost dictionary, the minimal opaque cost of the next sealed layer is the size of
its active primitive footprint:
\[
  \mu_{n+1}=|\mathrm{Foot}^{(d)}_{n,\mathrm{full}}|.
\]
Full coupling identifies this footprint with the whole historical basis window.  Lemma~\ref{lem:window-cardinality} gives
\[
  |\mathrm{Foot}^{(d)}_{n,\mathrm{full}}|
  =
  \sum_{j=0}^{d-1} (\kappa_{n-j}+\mu_{n-j}).
\]
Substitution yields the displayed recurrence.
\end{proof}

The theorem is deliberately modest.  It does not prove that a system has window size \(d\),
and it does not prove that a declaration is fully coupled.  Those are semantic and
presentation-level hypotheses.  The theorem says that once those hypotheses are supplied, the
trace count is forced.

\begin{corollary}[Depth-two full-coupling recurrence]
\label{cor:depth-two-full-coupling-recurrence}
In the fixed fully coupled cubical sealed-extension calculus, the two-layer chronological
factorization theorem of Section~5 gives
\[
  \mu_{n+1}
  =
  \mu_n+\mu_{n-1}+\kappa_n+\kappa_{n-1}
  \qquad (n\geq 2).
\]
\end{corollary}

\begin{proof}
The fixed cubical calculus has chronological window size two for primitive public support, by
the recent-history factorization theorem.  Apply
Theorem~\ref{thm:full-coupling-affine-recurrence} with \(d=2\).
\end{proof}

Let
\[
  \tau_n := \sum_{i=1}^{n}\mu_i
\]
be the cumulative minimal opaque trace cost through stage \(n\).  This cumulative quantity is
not used in the semantic depth theorem; it is useful only for comparing growth rates across
window sizes.

\subsection{The depth-one contrast}
\label{subsec:depth-one-growth}

The recurrence also explains why depth one is a genuinely different scaling regime.  If all
primitive support factors through only the immediately preceding layer, then the next cost is
linear in the previous public basis size, not in two previous basis sizes.

\begin{corollary}[Affine one-step growth at depth one]
\label{cor:depth-one-affine-growth}
Assume obligation depth one is realized by a one-layer chronological window, and assume every
sealed layer contributes the same core payload size \(c\).  Then
\[
  \mu_{n+1}=\mu_n+c.
\]
If the initial state exports no prior sealed layers, then
\[
  \mu_1=0,
  \qquad
  \mu_n=(n-1)c,
\]
and therefore
\[
  \tau_n = c\,\frac{n(n-1)}{2}.
\]
Thus the layerwise minimal opaque cost grows linearly, while the cumulative opaque cost grows
quadratically.
\end{corollary}

\begin{proof}
Substitute \(d=1\) into Theorem~\ref{thm:full-coupling-affine-recurrence}.  Only the
\(j=0\) term remains, so
\[
  \mu_{n+1}=\mu_n+\kappa_n.
\]
Under the uniform-payload assumption \(\kappa_n=c\), this is
\(\mu_{n+1}=\mu_n+c\).  If the initial state has no sealed historical interface, then the
first layer has no prior basis against which to pay trace, so \(\mu_1=0\).  Induction gives
\(\mu_n=(n-1)c\), and summing the arithmetic progression gives the formula for \(\tau_n\).
\end{proof}

This is a useful sanity check.  In an extensional or UIP-like regime where comparison data
collapses to boundary action, binary trace is not a genuine primitive public obligation.  The
scaling law reflects that collapse.  The fixed cubical theorem is different because univalence
supplies a real binary obstruction.

\subsection{The depth-two Fibonacci endpoint}
\label{subsec:fibonacci-endpoint}

We now specialize the depth-two recurrence to constant payload size.  This is the place where
Fibonacci growth appears.  It should be read as an endpoint of the full-coupling regime, not
as the definition of coherence depth two.

Let \((F_n)_{n\geq 1}\) be the Fibonacci sequence with
\[
  F_1=F_2=1,
  \qquad
  F_{n+1}=F_n+F_{n-1}.
\]

\begin{corollary}[Affine Fibonacci envelope]
\label{cor:affine-fibonacci-envelope}
In the fixed fully coupled cubical sealed-extension calculus, assume obligation depth two is
realized by the two-layer chronological window and every sealed layer contributes the same
core payload size \(c\).  Then
\[
  \mu_{n+1}=\mu_n+\mu_{n-1}+2c.
\]
Define the shifted sequence
\[
  U_n := \mu_n+2c.
\]
Then
\[
  U_{n+1}=U_n+U_{n-1}.
\]
If the initial state exports no prior sealed layers, then
\[
  \mu_1=0,
  \qquad
  \mu_2=c,
\]
and consequently
\[
  U_n=cF_{n+2},
  \qquad
  \mu_n=c(F_{n+2}-2),
\]
while the cumulative minimal opaque cost satisfies
\[
  \tau_n=c(F_{n+4}-2n-3).
\]
\end{corollary}

\begin{proof}
Corollary~\ref{cor:depth-two-full-coupling-recurrence} gives
\[
  \mu_{n+1}
  =
  \mu_n+\mu_{n-1}+\kappa_n+\kappa_{n-1}.
\]
Under the uniform-payload assumption \(\kappa_n=c\), this becomes
\[
  \mu_{n+1}=\mu_n+\mu_{n-1}+2c.
\]
Adding \(2c\) to both sides yields
\[
  \mu_{n+1}+2c
  =
  (\mu_n+2c)+(\mu_{n-1}+2c),
\]
which is exactly \(U_{n+1}=U_n+U_{n-1}\).

For the bootstrap values, the first layer sees no previous sealed public interface, so
\(\mu_1=0\).  The second layer sees only the first layer's payload basis, of size \(c\), and
therefore \(\mu_2=c\).  Hence
\[
  U_1=2c=cF_3,
  \qquad
  U_2=3c=cF_4.
\]
Since \((U_n)\) and \((cF_{n+2})\) satisfy the same recurrence and agree at the first two
stages, induction gives \(U_n=cF_{n+2}\).  Subtracting \(2c\) gives the formula for
\(\mu_n\).  Finally,
\[
  \tau_n
  =
  \sum_{i=1}^{n} c(F_{i+2}-2)
  =
  c(F_{n+4}-2n-3),
\]
using the standard identity \(\sum_{i=1}^n F_{i+2}=F_{n+4}-3\).
\end{proof}

The first few layers are shown below.  The table uses the bundled per-site convention: one
active primitive site contributes one public trace schema for the whole payload package.

\begin{center}
\begin{tabular}{c c c c}
\hline
\(n\) & \(\kappa_n\) & \(\mu_n\) & \(|S_{\mathrm{bas}}(L_n)|=\kappa_n+\mu_n\) \\
\hline
1 & \(c\) & \(0\) & \(c\) \\
2 & \(c\) & \(c\) & \(2c\) \\
3 & \(c\) & \(3c\) & \(4c\) \\
4 & \(c\) & \(6c\) & \(7c\) \\
5 & \(c\) & \(11c\) & \(12c\) \\
\hline
\end{tabular}
\end{center}

\begin{corollary}[Asymptotic inflation factor]
\label{cor:asymptotic-inflation-factor}
Under the bootstrap and uniform-payload hypotheses of
Corollary~\ref{cor:affine-fibonacci-envelope}, the structural inflation factor
\[
  \Phi_n := \frac{\mu_n}{\mu_{n-1}}
\]
converges, for \(n\) large enough that the ratio is defined, to the Golden Ratio
\[
  \varphi=\frac{1+\sqrt{5}}{2}.
\]
\end{corollary}

\begin{proof}
By Corollary~\ref{cor:affine-fibonacci-envelope},
\[
  \frac{\mu_n}{\mu_{n-1}}
  =
  \frac{F_{n+2}-2}{F_{n+1}-2}.
\]
The additive constants are asymptotically negligible relative to the Fibonacci terms, and the
standard ratio convergence \(F_{n+2}/F_{n+1}\to\varphi\) gives the result.
\end{proof}

This corollary is intentionally placed after all the hypotheses have been stated.  The Golden
Ratio appears only after depth two, two-layer chronology, full coupling, bundled per-site
counting, uniform payload, and the empty-history bootstrap have all been imposed.

\subsection{Robustness under refactoring}
\label{subsec:scaling-robustness}

The preceding counts are not source-level line counts.  They are cardinalities of normalized
public basis sets.  This makes the recurrence stable under the harmless refactorings that are
common in type-theoretic presentations.

\begin{corollary}[Robustness under admissible refactoring]
\label{cor:robustness-under-refactoring}
Let two presentations of the same sealed payload or structural obligation be related by a
canonical telescope isomorphism: rebundling or splitting dependent records, reassociation of
Sigma-types, currying or uncurrying iterated function arguments, renaming of bound variables,
or transparent internal normalization that does not change which fields are primitively
inhabited.  Then the induced sets of atomic payload schemas, primitive trace schemas, and
historical supports are canonically bijective.

Consequently, the quantities
\[
  \kappa(X),
  \qquad
  \mu(e),
  \qquad
  |S_{\mathrm{bas}}(L_n)|,
  \qquad
  |\mathrm{Foot}^{(d)}_n|
\]
are unchanged by such refactorings.
\end{corollary}

\begin{proof}
The listed refactorings preserve the same dependent typing judgment and differ only by
canonical rebracketing, rebundling, or transparent elaboration of its telescope.  After
flattening the Sigma-spine and discarding uniformly synthesized fields, both presentations
reduce to the same counting normal form.  Their atomic payload generators, primitive trace
fields, and supporting layer indices therefore correspond bijectively.  Since the metrics used
in this section are cardinalities of those normal-form sets, they are invariant.
\end{proof}

This robustness is why the recurrence is stated in terms of \(S_{\mathrm{bas}}(L_n)\) rather
than in terms of concrete syntax.  A source language may expose a single record with many
fields or many small declarations; the theorem counts the normalized active basis after the
adequacy bridge has decided which fields are primitive and which are derived.

\subsection{Abstract wrappers and scope of generalization}
\label{subsec:abstract-wrappers}

The fixed cubical theorem and the counting theorem now separate cleanly.  The counting theorem
needs only a window and a full-coupling footprint.  The stronger depth theorem needs a
horn-computational bridge: higher structural obligations must be represented as theorem-facing
open-box extension packages, those packages must be contractible, and explicit higher fields
must be replaced by cubical terms before \(\mu\) is computed.

This subsection packages that distinction.  The first theorem is the stronger transfer
criterion.  It is the right formulation of the broad guiding conjecture: a foundation or
surface calculus inherits the depth-two theorem only after supplying these clauses.  The second
wrapper is weaker and retains only the exact-depth and window information needed for the
affine recurrence.

\begin{theorem}[Abstract horn-computational foundation interface]
\label{thm:abstract-horn-computational-interface}
Let a sealed extension calculus satisfy the following package.  The fiber \(H_k(\partial)\)
below is not the ambient type of all semantic fillers for a closed boundary.  It is the
theorem-facing exported horn-extension package selected by the sealed calculus after
normalization.

\begin{enumerate}
\item \textbf{Horn representation.}  For every candidate \(X\), every \(k\geq 3\), and every
boundary \(\partial : \mathrm{Real}_{k-1}(X)\), the depth-\(k\) obligation over \(\partial\) is
canonically represented by a horn-extension fiber
\[
  H_k(\partial)
  :=
  \sum_{\gamma:\Gamma^{(k)}_{\mathrm{miss}}(\partial)}
  \mathrm{Fill}^{(k)}_{\partial}(\gamma),
\]
and there is a canonical equivalence
\[
  \mathrm{Real}_k(X)
  \simeq
  \sum_{\partial:\mathrm{Real}_{k-1}(X)} H_k(\partial).
\]

\item \textbf{Theorem-facing contractibility.}  Every horn-extension fiber
\(H_k(\partial)\) appearing in the previous clause is canonically contractible.

\item \textbf{Canonical public presentations.}  Every admissible surface extension
declaration admits a canonical normalized trace presentation, with presentation equivalence
and a primitive-versus-derived classification of trace fields.

\item \textbf{Computational replacement.}  In the canonical presentation, every historical
trace field of arity \(k\geq 3\) corresponding to a horn factor of the first clause is
transparently derived.  Replacing it by the canonically synthesized horn witness preserves the
canonical normalized trace presentation.

\item \textbf{Two-layer primitive support.}  Every primitive field in a canonical normalized
trace presentation depends only on the last two layers, up to the factorization relation used
for normalized public signatures.
\end{enumerate}

Then the calculus satisfies all three conclusions below.
\begin{enumerate}
\item For every candidate \(X\) and every \(k\geq 2\), there is a canonical equivalence
\[
  \mathrm{Real}_k(X) \simeq \mathrm{Real}_2(X).
\]
\item Every admissible surface declaration is presentation-equivalent to one whose normalized
trace signature contains no irreducible primitive historical field of arity \(k\geq 3\).
\item Primitive historical support factors through a chronological window of size two.
\end{enumerate}
In particular, the exact obligation-depth bound \(d_{\mathrm{obl}}\leq 2\) and the
minimal-signature bound \(d_\mu\leq 2\) both hold.
\end{theorem}

\begin{proof}
By horn representation, the step from depth \(k-1\) to depth \(k\) for \(k\geq 3\) decomposes
as a boundary in \(\mathrm{Real}_{k-1}(X)\) together with one theorem-facing horn factor
\(H_k(\partial)\).  By theorem-facing contractibility, projection to the boundary is a
canonical equivalence.  Iterating this projection from depth \(k\) down to depth two gives
\[
  \mathrm{Real}_k(X) \simeq \mathrm{Real}_2(X).
\]
This proves exact stabilization above binary arity.

For the public-signature conclusion, take an admissible surface declaration and pass to its
canonical normalized trace presentation.  Any historical field of arity \(k\geq 3\) that
corresponds to one of the horn factors is, by computational replacement, replaced by the
canonically synthesized horn witness without changing the canonical presentation.  Repeating
this finite replacement for all such fields gives a presentation-equivalent declaration whose
\(\mu\)-minimal signature contains no irreducible primitive field above binary arity.  The
primitive-support clause is exactly the chronological-window statement.
\end{proof}

\begin{corollary}[Fixed cubical instantiation]
\label{cor:fixed-cubical-instantiation}
The fixed cubical sealed-extension calculus over the chosen CCHM-style core satisfies the
horn-computational interface of
Theorem~\ref{thm:abstract-horn-computational-interface}.  Consequently the exact
stabilization theorem, the higher-field elimination theorem, and the recent-history
factorization theorem of Section~5 are the concrete instance of this abstract pattern proved
for the fixed calculus.
\end{corollary}

\begin{proof}
Horn representation is the structural horn-to-open-box reduction of Section~5.
Theorem-facing contractibility is the open-box extension theorem, whose center is computed by
cubical composition and filling.  Canonical public presentations and presentation equivalence
are supplied by the adequacy and normalization interface of Section~3.  Computational
replacement is the theorem that explicit higher structural fields elaborate to cubical horn
terms and disappear from \(\mu\)-minimal public trace.  Two-layer primitive support is the
recent-history factorization theorem.  Therefore the abstract interface applies to the fixed
cubical calculus.
\end{proof}

The next theorem records the weaker wrapper needed only for the recurrence.  It deliberately
forgets the canonical-presentation and computational-replacement clauses.  This is useful for
small toy disciplines that have an explicit primitive binary layer and a stipulated two-layer
window, even if they are not claimed to instantiate the full horn-computational cubical bridge.

\begin{theorem}[Depth-two law for a specified two-dimensional discipline]
\label{thm:specified-2d-discipline}
Let a fully coupled foundation or extension discipline satisfy the following two properties.
\begin{enumerate}
\item There is a candidate \(X\) whose realized binary obligation object is not canonically
equivalent to its unary obligation object:
\[
  \mathrm{Real}_2(X) \not\simeq \mathrm{Real}_1(X).
\]
\item For every candidate \(X\) and every \(k\geq 2\), there is a canonical equivalence
\[
  \mathrm{Real}_k(X) \simeq \mathrm{Real}_2(X),
\]
and every surviving binary clause factors through a chronological window of size two.
\end{enumerate}
Then the foundation has obligation depth two.  Under the uniform payload hypothesis
\(\kappa_n=c\), its fully coupled minimal opaque cost satisfies
\[
  \mu_{n+1}=\mu_n+\mu_{n-1}+2c,
\]
and the shifted sequence \(U_n:=\mu_n+2c\) is Fibonacci.
\end{theorem}

\begin{proof}
The second property gives exact stabilization at depth two and a genuine two-layer
chronological window.  The recurrence follows from
Theorem~\ref{thm:full-coupling-affine-recurrence} with \(d=2\) and \(\kappa_n=c\).  The
first property excludes depth one.  Therefore the least stabilizing obligation depth is exactly
two, and the Fibonacci shift is the one proved in
Corollary~\ref{cor:affine-fibonacci-envelope}.
\end{proof}

\begin{corollary}[Concrete strict/fibrant two-level toy instance]
\label{cor:concrete-2ltt-instance}
There is a concrete strict/fibrant two-level type-theoretic toy discipline satisfying the
hypotheses of Theorem~\ref{thm:specified-2d-discipline}.  In this discipline, arities zero and
one carry no primitive obligations, arity two carries a primitive fibrant binary comparison,
every arity \(k\geq 3\) carries the same comparison schema only as transparently derived data,
and every surviving primitive support factors through the previous two layers.  Consequently,
this discipline has obligation depth two and chronological window size two.
\end{corollary}

\begin{proof}
Take the described strict/fibrant discipline as the model.  Since the binary comparison is
primitive while the unary layer is empty, the lower-bound property holds.  Since every arity
\(k\geq 2\) uses the same binary comparison schema, exact stabilization above arity two is
immediate.  Since primitive support is stipulated to use only the previous two layers, the
chronological-window property also holds.  Theorem~\ref{thm:specified-2d-discipline} applies.
\end{proof}

\begin{remark}[Why plain Book HoTT is not claimed]
\label{rem:plain-hott-not-claimed}
The identity-elimination principle of Book HoTT reduces families over a variable path to the
reflexivity case, but it does not by itself provide the cubical open-box computations used in
the fixed calculus.  In particular, the structural horn-extension problems of Section~5 do not
automatically become theorem-facing contractible open boxes in plain Book HoTT.  Therefore
this paper does not claim that plain HoTT instantiates the horn-computational interface of
Theorem~\ref{thm:abstract-horn-computational-interface}, nor even that it satisfies the
weaker exact/window package of Theorem~\ref{thm:specified-2d-discipline}.  Plain HoTT remains
outside the theorem-level instantiations considered here.
\end{remark}

The scope of the generalization claims can be summarized as follows.

\begin{center}
\scriptsize
\setlength{\tabcolsep}{2pt}
\begin{tabular}{p{0.22\linewidth}p{0.17\linewidth}p{0.22\linewidth}p{0.17\linewidth}p{0.07\linewidth}}
\hline
Foundation or discipline & Binary lower bound & Exact stabilization / window & Bridge to \(\mu\) & Claimed? \\
\hline
Chosen CCHM-style cubical core / fixed raw calculus
& yes; explicit swap obstruction
& yes; horn reduction, open-box stabilization, and recent-history factorization
& yes; canonical presentation and computational replacement
& yes \\
\hline
Other cubical variants
& only if a binary witness is checked case by case
& only if the horn-representation, theorem-facing fiber, and window clauses are verified
& only if the canonical-trace bridge clauses are verified
& no \\
\hline
Explicit strict/fibrant two-level toy discipline
& yes; primitive binary witness
& yes; exact stabilization and size-two window by construction
& not supplied here; only the weaker wrapper is claimed
& yes, for that discipline \\
\hline
Plain Book HoTT
& no theorem isolated here
& no cubical-style exact-stabilization argument supplied here
& no
& no \\
\hline
\end{tabular}
\end{center}

This table is meant to keep the boundary honest.  The paper proves a fixed cubical theorem, a
counting consequence under full coupling, and a small two-dimensional toy wrapper.  Broader
transfer to other foundations is a program: each proposed target must instantiate the
horn-computational interface, or else explicitly work only with the weaker exact/window
package needed for arithmetic.
